% =====================================================================
%  Abstract – „Making-of“ und Reflexion (Phase 3)
%  Formalia laut Aufgabenstellung 4.3:
%   - Umfang: 2 Seiten TEXTTEIL (Titelblatt/Literaturverzeichnis zählen nicht)
%   - DIN A4, Ränder 2 cm, Arial 11 pt, Überschriften 12 pt, Blocksatz,
%     Silbentrennung, Zeilenabstand 1,5, Fußnoten Arial 10 pt,
%     6 pt Abstand nach Absätzen  -> Klassenoption „abstract“
%  Inhalt (Aufgabenstellung 1.3 / Prüfungsleitfaden Phase 3):
%   - Wie verlief die Entwicklung? Was funktionierte gut?
%   - Was könnte beim nächsten Projekt besser gemacht werden?
%   - Welche wichtigen Erkenntnisse wurden gewonnen?
%   - Umgang mit unerwartet aufgetretenen Problemen
%   - Einhaltung des ursprünglich anvisierten Zeitplans
%   - Kurze Vorstellung von Zielsetzung/Idee/Konzept/Methodik;
%     Ergebnis vs. Ziel – falls nicht erreicht: warum nicht?
%  Selbstreflexiver Text: „ich“ ist hier zulässig (Richtlinien Kap. 1).
%  Dateiname der Abgabe: Nachname-Vorname_Matrikelnummer_PSE_P3_Abstract.pdf
%  tools/build_docs.py prüft, dass der Textteil genau 2 Seiten umfasst.
% =====================================================================
\documentclass[abstract]{iu-pse}
\input{meta}
\addbibresource{literatur.bib}

\dokumenttitel{Abstract}
\dokumentuntertitel{\Projektname{} -- Making-of und Reflexion}
\abgabedatum{\AbgabePhaseDrei}

\begin{document}
\titelblatt
\hauptteil

\section{Zielsetzung und Konzept}
\TODO{Kurze Vorstellung von Ziel, Idee, Konzept und Methodik}

\section{Verlauf der Entwicklung}
\TODO{Making-of: Phasen, wichtige Entscheidungen, eingesetzte Ressourcen und Software}

\section{Probleme und Lösungen}
\TODO{Unerwartete Probleme und wie damit umgegangen wurde}

\section{Zeitplan und Zielerreichung}
\TODO{Soll-Ist-Vergleich des Zeitplans; Entspricht das Ergebnis dem Ziel? Falls nein, warum nicht?}

\section{Erkenntnisse und Ausblick}
\TODO{Was lief gut, was beim nächsten Projekt besser machen, wichtigste Lessons Learned}

% Nur falls im Abstract zitiert wird:
% \literaturverzeichnis
\end{document}
